\documentclass[10pt, aspectratio=169]{beamer}
\usepackage{amsmath, amssymb, amsthm}
\usepackage{xcolor}

\usetheme{Madrid}
\usecolortheme{whale}

\title[Phase 2: The Integer Engine]{Phase 2: The Integer Engine and The Noise Crisis}
\subtitle{Fully Homomorphic Encryption over the Integers}
\author{Paper Presentation}
\date{\today}

\begin{document}

\frame{\titlepage}

\begin{frame}{Outline: Phase 2}
    \tableofcontents
\end{frame}

\section{The Symmetric \& Public-Key SHE Schemes}

\begin{frame}{Building the Engine: The Symmetric "Toy" Scheme}
    Before we look at the complex public-key version, we start with a beautifully simple symmetric-key foundation.
    \pause
    
    \vspace{0.3cm}
    \textbf{Parameters:} Let $\eta$ be the bit-length of the secret key, and $\rho$ be the bit-length of the noise.
    \pause
    
    \begin{block}{The Symmetric Algorithms}
    \begin{itemize}
        \item \textbf{KeyGen:} The secret key $p$ is a randomly chosen odd integer $p \in [2^{\eta-1}, 2^{\eta})$.
        \item \textbf{Encrypt($p, m$):} To encrypt a bit $m \in \{0,1\}$, choose a large random integer $q$ and a small random noise integer $r$ (where $|r| < 2^{\rho}$). 
        $$c = p \cdot q + 2r + m$$
        \item \textbf{Decrypt($p, c$):} Output $m' = (c \pmod p) \pmod 2$.
    \end{itemize}
    \end{block}
\end{frame}

\begin{frame}{Intuition: Why Symmetric Decryption Works}
    Let's break down the decryption equation: $m' = (c \pmod p) \pmod 2$.
    \pause
    
    \vspace{0.3cm}
    \begin{enumerate}
        \item Substitute $c$: $(p \cdot q + 2r + m \pmod p)$
        \pause
        \item Because $p \cdot q$ is an exact multiple of $p$, it vanishes perfectly. We are left with $(2r + m \pmod p)$.
        \pause
        \item \textbf{The Critical Condition:} The modulo $p$ operation only leaves $2r+m$ unchanged if its absolute value is strictly less than half the modulus. We strictly require:
        $$|2r + m| < p/2$$
        \pause
        \item Assuming the noise is small enough, we take $(2r + m) \pmod 2$. The $2r$ is even, so it vanishes, leaving exactly $m$.
    \end{enumerate}
    \pause
    
    \begin{block}{Motivation}
    To anyone without $p$, the ciphertext $c$ just looks like a massive, random integer. But algebraically, it is a point "tethered" to a multiple of $p$ by a short, noisy string ($2r+m$).
    \end{block}
\end{frame}

\begin{frame}{Transitioning to Public-Key}
    A true FHE scheme requires public-key infrastructure. We achieve this by publishing a massive list of "encryptions of zero."
    \pause

    \vspace{0.2cm}
    \textbf{New Parameters (based on security parameter $\lambda$):}
    \begin{itemize}
        \item $\gamma$: Bit-length of the integers in the public key.
        \item $\tau$: Number of integers in the public key.
        \item $\rho'$: Secondary noise parameter for encryption $\rho' = \rho + \omega(\log \lambda)$.
    \end{itemize}
    \pause

    \vspace{0.2cm}
    \begin{block}{The Noise Distribution $\mathcal{D}_{\gamma,\rho}(p)$}
    To generate an "encryption of zero", choose $q \leftarrow \mathbb{Z} \cap [0, 2^\gamma/p)$ and $r \leftarrow \mathbb{Z} \cap (-2^\rho, 2^\rho)$. 
    Output $x = pq + r$.
    \end{block}
\end{frame}

\begin{frame}{The Public-Key Algorithms}
    \textbf{KeyGen:}
    \begin{itemize}
        \item Sample $\tau + 1$ integers $x_0, x_1, \dots, x_\tau$ from $\mathcal{D}_{\gamma,\rho}(p)$. 
        \item Relabel so $x_0$ is the largest. Restart unless $x_0$ is odd and its noise $r_p(x_0)$ is even.
        \item Public key $pk = \langle x_0, x_1, \dots, x_\tau \rangle$.
    \end{itemize}
    \pause
    
    \vspace{0.3cm}
    \textbf{Encrypt($pk, m$):}
    Choose a random subset $S \subseteq \{1, 2, \dots, \tau\}$ and a noise integer $r \in (-2^{\rho'}, 2^{\rho'})$.
    $$c = \left[ m + 2r + 2\sum_{i \in S} x_i \right]_{x_0}$$
    \pause
    
    \begin{block}{Intuition}
    Every $x_i$ is effectively $0$ (plus noise). Summing a random subset of $x_i$ computes $0+0+0\dots$, creating a fresh, randomized encryption of zero. We multiply by 2 to keep the noise even, add our own fresh noise $2r$, and add $m$. The final modulo $x_0$ prevents the physical bit-size from growing infinitely large.
    \end{block}
\end{frame}

\section{Rigorous Proof of Correctness}

\begin{frame}{Lemma A.1: The Structure of a Fresh Ciphertext}
    To prove this scheme is homomorphic, we first mathematically prove that a fresh public-key ciphertext has the exact same structure as our symmetric toy scheme.
    \pause
    
    \begin{lemma}[A.1]
    Let $c \leftarrow \text{Encrypt}(pk, m)$. Then, $c = a\cdot p + (2b + m)$ for some integers $a$ and $b$ where the noise is bounded by $|2b+m| \le \tau 2^{\rho+2}$.
    \end{lemma}
    \pause
    
    \vspace{0.3cm}
    \textbf{Proof Sketch:}
    \begin{itemize}
        \item By definition, $c = m + 2r + 2\sum_{i \in S} x_i + k \cdot x_0$.
        \item We know every $x_i = q_i p + r_i$. Substitute this in:
        $$c = p \cdot \left(kq_0 + 2\sum_{i \in S} q_i\right) + \left(m + 2r + 2k \cdot r_0 + 2\sum_{i \in S} r_i\right)$$
        \item The first term is a multiple of $p$ (our $a \cdot p$). The second term is our accumulated noise ($2b+m$). Its parity is $m$, and its magnitude is bounded.
    \end{itemize}
\end{frame}

\begin{frame}{Homomorphic Evaluation: Addition}
    Let the Evaluate algorithm simply be integer addition and multiplication. 
    Let $c_1 = q_1 p + 2r_1 + m_1$ and $c_2 = q_2 p + 2r_2 + m_2$.
    \pause

    \vspace{0.3cm}
    \textbf{Homomorphic Addition:}
    $$c_1 + c_2 = p(q_1 + q_2) + 2(r_1 + r_2) + (m_1 + m_2)$$
    \pause
    
    \vspace{0.3cm}
    \textbf{Observation:}
    The structural format $a\cdot p + \text{noise}$ is perfectly preserved! We have a new multiple of $p$, plus $2 \times (\text{new noise})$, plus the sum of the messages. 
    
    \vspace{0.2cm}
    \textbf{Crucially:} The noise grows \textbf{linearly} ($r_1 + r_2$).
\end{frame}

\begin{frame}{Homomorphic Evaluation: Multiplication}
    Now, let us multiply the ciphertexts:
    $$c_1 \cdot c_2 = (q_1 p + 2r_1 + m_1)(q_2 p + 2r_2 + m_2)$$
    \pause
    
    Expanding this algebraically yields:
    $$c_1 \cdot c_2 = p \big( p \cdot q_1 q_2 + q_1(2r_2 + m_2) + q_2(2r_1 + m_1) \big)$$
    $$+ \ 2 \big( 2r_1 r_2 + r_1 m_2 + r_2 m_1 \big) + m_1 m_2$$
    \pause
    
    \vspace{0.3cm}
    \textbf{Observation:}
    The structure is still preserved. We have a massive multiple of $p$ (first term), plus $2 \times (\text{mixed noise})$, plus the product of the messages $m_1 m_2$.
    
    \vspace{0.2cm}
    \textbf{The Crisis:} Look at the dominant noise term: $\mathbf{2r_1 r_2}$. When we multiply ciphertexts, the noise does not add—it \textbf{multiplies}.
\end{frame}

\section{The Multiplicative Bottleneck}

\begin{frame}{The Decryption Constraint \& The Noise Crisis}
    \textbf{Lemma A.2 Constraint:} 
    Recall our strict rule for correct decryption. The modulo $p$ operation $(c \pmod p)$ will completely destroy the plaintext unless the absolute value of the total accumulated noise is strictly less than $p/2$.
    \pause
    
    \vspace{0.3cm}
    \begin{block}{The Crisis of the Square}
    Because multiplication squares the magnitude of the noise, the noise bit-length doubles with every sequential multiplication. It grows exponentially with the multiplicative depth of the circuit. We will rapidly hit the $p/2$ boundary.
    \end{block}
\end{frame}

\begin{frame}{Formalizing the Bottleneck: Lemma 3.5}
    Exactly how deep can our circuit be before the noise destroys the data?
    \pause
    
    \begin{lemma}[3.5: Permitted Polynomials]
    Let $f(x_1, \dots, x_t)$ be the multivariate polynomial computed by a circuit, and let $d$ be its degree. The scheme can correctly evaluate this polynomial as long as:
    $$d \le \frac{\eta - 4 - \log||\vec{f}||_1}{\rho' + 2}$$
    where $||\vec{f}||_1$ is the $l_1$ norm of the coefficient vector of $f$.
    \end{lemma}
    \pause
    
    \begin{block}{Intuition: The Noise Budget}
    \begin{itemize}
        \item Think of $\eta$ (the bit-size of $p$) as our total \textbf{Noise Budget}.
        \item Think of $\rho'$ (the initial noise size) as the \textbf{cost} of one multiplication.
        \item Therefore, the maximum multiplicative depth (degree $d$) we can handle is roughly our total budget divided by the cost per operation: $\approx \eta / \rho'$.
    \end{itemize}
    \end{block}
\end{frame}

\begin{frame}{Conclusion of Phase 2}
    \textbf{Where do we stand?}
    \pause
    
    \begin{enumerate}
        \item We have successfully built an encryption scheme that can evaluate additions and multiplications without the server knowing the underlying data.
        \pause
        \item However, because of the multiplicative noise bottleneck, this scheme is strictly a \textbf{Somewhat Homomorphic Encryption (SHE)} scheme.
        \pause
        \item If we evaluate a circuit deeper than $d \approx \eta / \rho'$, the noise wraps around the modulus $p$, and the data is lost forever.
    \end{enumerate}
    \pause
    
    \vspace{0.5cm}
    \begin{center}
        \textit{To achieve our dream of infinite FHE computation, we must find a way to evaluate our own decryption circuit within this strict noise budget. This leads us to Phase 3...}
    \end{center}
\end{frame}

\section{References}
\begin{frame}{References}
    \begin{thebibliography}{9}
        \bibitem{DGHV10}
        Marten van Dijk, Craig Gentry, Shai Halevi, and Vinod Vaikuntanathan.
        \textit{Fully Homomorphic Encryption over the Integers}. 
        Eurocrypt, 2010.
    \end{thebibliography}
\end{frame}

\end{document}